% BHU PYQ -- Manually transcribed & error-corrected
% Paper: MTB-101 : Calculus I
% Exam: B.Sc. (Hons.) Semester I Examination, 2016-17

\documentclass[11pt,a4paper]{article}
\usepackage[a4paper,top=2.5cm,bottom=2.5cm,left=2.8cm,right=2.8cm,headheight=14pt]{geometry}
\usepackage{amsmath,amssymb}
\usepackage[T1]{fontenc}
\usepackage[utf8]{inputenc}
\usepackage{lmodern}
\usepackage{microtype}
\usepackage{fancyhdr}
\usepackage{enumitem}
\usepackage{hyperref}
\usepackage{lastpage}
\usepackage{parskip}

\hypersetup{colorlinks=true,urlcolor=black,linkcolor=black,
  pdftitle={MTB-101: Calculus I -- BHU B.Sc. Sem I 2016-17}}
\pagestyle{fancy}\fancyhf{}
\renewcommand{\headrulewidth}{0.4pt}\renewcommand{\footrulewidth}{0.4pt}
\fancyhead[L]{\small   \textit{Banaras Hindu University}}
\fancyhead[R]{\small   \textit{MTB-101: Calculus I}}
\fancyfoot[C]{\small Page  \thepage\ of \pageref{LastPage}}


\newlist{parts}{enumerate}{2}
\setlist[parts,1]{label=(\alph*),leftmargin=*,itemsep=5pt,topsep=3pt}
\setlist[parts,2]{label=(\roman*),leftmargin=*,itemsep=3pt,topsep=2pt}

\newcommand{\pts}[1]{\hfill   \textit{[#1]}}


\begin{document}

\begin{center}
  {\large\bfseries BANARAS HINDU UNIVERSITY}\\[2pt]
  {\normalsize B.Sc. (Hons.) --- Semester I Examination, 2016-17}\\[6pt]
  \rule{0.8\linewidth}{0.5pt}\\[6pt]
  {\large\bfseries Paper No. MTB-101: Calculus I}\\[4pt]
  {\normalsize Time: 3 Hours \hfill Full Marks: 70}
\end{center}
\rule{\linewidth}{0.4pt}
\smallskip



\noindent   \textit{Instructions: Answer   \textbf{five} questions including Question~1 which is compulsory. Figures in right-hand margin indicate marks.}
\bigskip




\noindent   \textbf{Question 1.}   \textit{(Compulsory --- 2 marks each.)}

\begin{parts}
  \item Show that $\lim_{n  o \infty} \frac{1 + 3 + 5 + \dots + (2n-1)}{n^2} = 1$.
  \item Show that the series $\sum_{n=1}^\infty \frac{1}{\sqrt{n}}$ is not convergent.
  \item Test the continuity of the function defined by
  \[
  f(x) = 
\begin{cases} \frac{|x-a|}{x-a}, & x 
eq a \\ 1, & x = a \end{cases}
  \]
  at $x = a$. If discontinuous, write the type of discontinuity at $x = a$.
  \item Find the value of the derivative of the function $f(x) = |x-1| + |x-3|$ at $x=2$.
  \item If $y = \cos^2 x$, then find $y_n$.
  \item If $e^x = a_0 + a_1 x + a_2 x^2 + a_3 x^3 + \dots$, then find the value of $a_3$.
  \item By using the definition of limit, show that $\lim_{x o 2} \frac{1}{(x-2)^2}  o \infty$.
  \item Discuss the applicability of Rolle's theorem to the function:
  \[
  f(x) = 
\begin{cases} x^2+1, & 0 \le x < 1 \\ 3-x, & 1 \le x \le 2 \end{cases}
  \]
  \item Find all the asymptotes of the curve $y^2(x^2 - y^2) = a^2(x^2 + y^2)$, $a > 0$.
  \item Find the nature of the multiple point $(0, 0)$ on the curve $x^2(x-y) + y^3 = 0$.
  \item By using definite integration, find the limit:
  \[
  \lim_{n o\infty} \left[ \frac{\sqrt{n^2-1^2}}{n^2} + \frac{\sqrt{n^2-2^2}}{n^2} + \dots + \frac{\sqrt{n^2-(n-1)^2}}{n^2} \right]
  \]
\end{parts}

\medskip\rule{\linewidth}{0.2pt}

\medskip


\noindent   \textbf{Question 2.}

\begin{parts}
  \item Define the convergence of a sequence $(x_n)$. By using the definition, show that the sequence $(x_n)$ defined by $x_n = \sqrt{n+1} - \sqrt{n}$ is convergent. \pts{4}
  \item Test the convergence of the series:
  \[
  \frac{1}{2} + \frac{1 \cdot 2}{3 \cdot 5} + \frac{1 \cdot 2 \cdot 3}{3 \cdot 5 \cdot 7} + \frac{1 \cdot 2 \cdot 3 \cdot 4}{3 \cdot 5 \cdot 7 \cdot 9} + \dots
  \] \pts{4}
  \item Test the convergence of the series:
  \[
  \frac{1}{2} + \frac{2}{5} + \frac{3}{10} + \dots + \frac{n}{n^2+1} + \dots
  \] \pts{4}
\end{parts}

\medskip\rule{\linewidth}{0.2pt}

\medskip


\noindent   \textbf{Question 3.}

\begin{parts}
  \item Prove that every monotonic increasing sequence which is bounded above is convergent. \pts{6}
  \item If $f(x) = 
\begin{cases} x^2 \sin(1/x), & x 
eq 0 \\ 0, & x = 0 \end{cases}$, then show that the function is differentiable at $x = 0$ but the derivative is not continuous at $x = 0$. \pts{6}
\end{parts}

\medskip\rule{\linewidth}{0.2pt}

\medskip


\noindent   \textbf{Question 4.}

\begin{parts}
  \item Let $\lim_{x o a} f(x) = l$ and $\lim_{x o a} g(x) = m$. Show that $\lim_{x o a} f(x)g(x) = lm$. \pts{6}
  \item State and prove the intermediate value theorem for continuous functions. \pts{6}
\end{parts}

\medskip\rule{\linewidth}{0.2pt}

\medskip


\noindent   \textbf{Question 5.}

\begin{parts}
  \item State and prove Lagrange's mean value theorem. Provide its geometrical interpretation. \pts{6}
  \item If $y^{1/m} + y^{-1/m} = 2x$ for $m \in \mathbb{N}$, then prove that:
  \[
  (x^2-1) y_{n+2} + (2n+1) x y_{n+1} + (n^2-m^2) y_n = 0
  \] \pts{6}
\end{parts}

\medskip\rule{\linewidth}{0.2pt}

\medskip


\noindent   \textbf{Question 6.}

\begin{parts}
  \item Find all the asymptotes of the curve:
  \[
  x^3 - x^2y - xy^2 + y^3 + 2x^2 - 4y^2 + 2xy + x + y + 1 = 0
  \] \pts{6}
  \item State and prove Taylor's theorem in finite form with Lagrange's form of remainder. \pts{6}
\end{parts}

\medskip\rule{\linewidth}{0.2pt}

\medskip


\noindent   \textbf{Question 7.}

\begin{parts}
  \item Trace the curve:
  \[
  y^2(a+x) = x^2(3a-x)
  \] \pts{6}
  \item By applying the definition of definite integration as the limit of a sum, evaluate:
  \[
  \lim_{n o\infty} \left[ \left(1 + \frac{1}{n}\right)\left(1 + \frac{2}{n}\right)\dots\left(1 + \frac{n}{n}\right) \right]^{1/n}
  \] \pts{6}
\end{parts}

\vfill
\rule{\linewidth}{0.4pt}

\begin{center}\small
  \href{https://bhu-exam.vercel.app/}{Click here to visit our website}\\[4pt]
  \href{https://me-aryan.vercel.app/}{Click here to visit our contributor}\\[4pt]
  \href{https://quashan-me.vercel.app/}{Click here to visit our contributor}
\end{center}

\end{document}
